% Generated by Listas/make_statements.py from pratica-lecture-05.md, item 24. Do not edit.
\textbf{Source.} Gibbons (1992), \emph{Game Theory for Applied Economists}, §2.3 Problems, Problem 2.16, p.~136 {[}PDF 143{]}. Answer in the source: none in the book. Reproduced verbatim. (The one-shot model it builds on is Gibbons §2.1.C, pp.~64--68: the union sets the wage, then the firm sets employment.)



\begin{quote}

2.16. In the model of wages and employment analyzed in Section 2.1.C, the backwards-induction outcome is not socially efficient. In practice, however, a firm and a union negotiate today over the terms of a three-year contract, then negotiate three years from today over the terms of a second contract, and so on. Thus, the relationship may be more accurately characterized as a repeated game, as in Espinosa and Rhee (1989).



This problem derives conditions under which a subgame-perfect Nash equilibrium in the infinitely repeated game is Pareto-superior to the backwards-induction outcome of the one-shot game. Denote the union's utility and the firm's profit in the backwards-induction outcome of the one-shot game by \(U^*\) and \(\pi^*\), respectively. Consider an alternative utility-profit pair \((U, \pi)\) associated with an alternative wage-employment pair \((w, L)\). Suppose that the parties share the discount factor \(\delta\) (per three-year period). Derive conditions on \((w, L)\) such that (1) \((U, \pi)\) Pareto-dominates \((U^*, \pi^*)\) and (2) \((U, \pi)\) is the outcome of a subgame-perfect Nash equilibrium of the infinitely repeated game, where \((U^*, \pi^*)\) is played forever following a deviation.

\end{quote}



\textbf{Referred to in the statement.} The model of wages and employment of Section 2.1.C (Gibbons pp.~64--66 {[}PDF p.~71--73{]}) (reproduced from the book):



\begin{quote}

In Leontief's (1946) model of the relationship between a firm and a monopoly union (i.e., a union that is the monopoly seller of labor to the firm), the union has exclusive control over wages, but the firm has exclusive control over employment. (Similar qualitative conclusions emerge in a more realistic model in which the firm and the union bargain over wages but the firm retains exclusive control over employment.) The union's utility function is \(U(w, L)\), where \(w\) is the wage the union demands from the firm and \(L\) is employment. Assume that \(U(w, L)\) increases in both \(w\) and \(L\). The firm's profit function is \(\pi(w, L) = R(L) - wL\), where \(R(L)\) is the revenue the firm can earn if it employs \(L\) workers (and makes the associated production and product-market decisions optimally). Assume that \(R(L)\) is increasing and concave.



Suppose the timing of the game is: (1) the union makes a wage demand, \(w\); (2) the firm observes (and accepts) \(w\) and then chooses employment, \(L\); (3) payoffs are \(U(w, L)\) and \(\pi(w, L)\). {[}\ldots{]}



First, we can characterize the firm's best response in stage (2), \(L^*(w)\), to an arbitrary wage demand by the union in stage (1), \(w\). Given \(w\), the firm chooses \(L^*(w)\) to solve



\[\max_{L \ge 0} \pi(w, L) = \max_{L \ge 0} R(L) - wL,\]



the first-order condition for which is



\[R'(L) - w = 0.\]



{[}\ldots{]} To guarantee that the first-order condition \(R'(L) - w = 0\) has a solution, assume that \(R'(0) = \infty\) and that \(R'(\infty) = 0\) {[}\ldots{]}



We turn next to the union's problem at stage (1). Since the union can solve the firm's second-stage problem as well as the firm can solve it, the union should anticipate that the firm's reaction to the wage demand \(w\) will be to choose the employment level \(L^*(w)\). Thus, the union's problem at the first stage amounts to



\[\max_{w \ge 0} U(w, L^*(w)).\]

\end{quote}

